\begin{theorem}[DP33: optimal affine errors for the height-three dual]\label{DP33}
In the two-coordinate model of DP32 with $h=3$, write $X=X_1$, $Y=X_2$, $q=p^3$, $c_X=p/(q-1)$ and $c_Y=p^2/(q-1)$. For every integral $f$ the coefficientwise bounds improve to
\[
 v_p(D_X^nf)\ge\max\{0,\ceil{c_Xn-1}\},\qquad
 v_p(D_Y^nf)\ge\max\{0,\ceil{c_Yn-1}\}.
\]
Both slopes and both constant errors $1$ are optimal. Mixed derivatives have the sum of the two displayed integer exponents as a lower bound. In addition to the exact diagonal sequences in DP32,
\[
 v_p(D_X^{pq^r}Y(0))=\frac{p^2(q^r-1)}{q-1},\qquad
 v_p(D_Y^{p^2q^r}X(0))=\frac{p(q^{r+1}-1)}{q-1}.
\]
These statements include $p=2$.
\end{theorem}
\begin{proof}
The actual-dual model and its integral coordinates are already established by DP32. Its logarithm begins
\[
 \mathcal L_1=X+Y^{p^2}/p+X^q/p^2+\cdots,\qquad
 \mathcal L_2=Y+X^p/p+Y^q/p^2+\cdots.
\]
Let $\Delta_{a,b}$ be the integral coefficient operators of translation by $(S,T)$. They commute with one another and with $D_X,D_Y$, by commutativity of the group. In characteristic zero the translation operator is $\exp(\mathcal L_1(S,T)D_X+\mathcal L_2(S,T)D_Y)$. Comparing coefficients gives
\[
 \Delta_{p,0}=D_X^p/p!+D_Y/p,\qquad
 \Delta_{0,p^2}=D_Y^{p^2}/(p^2)!+D_X/p,\qquad
 \Delta_{0,p}=D_Y^p/p!.
\]
Consequently, with integral commuting operators $T,S$ and units $a,b$,
\begin{equation}\label{eq:dual-low-ops}
 D_X^p=aD_Y+pT,\quad a=-(p-1)!,\qquad
 D_Y^{p^2}=p^p bD_X+p^{p+1}S,\quad b=-\frac{(p^2)!}{p^{p+1}}.
\end{equation}
Also $D_Y^p$ is divisible by $p$ as an operator. Raising the first identity to $p$ shows $D_X^{p^2}$ is divisible by $p$: the leading term uses $D_Y^p$, and every remaining binomial term has a factor $p$. Thus for all $n$ the elementary bounds are $\floor{n/p^2}$ in direction $X$ and $\floor{n/p}$ in direction $Y$.

Raise the first identity in \eqref{eq:dual-low-ops} to $p^2$. Its nonleading term indexed by $1\le j\le p^2$ has valuation at least
\[
 v_p\binom{p^2}{j}+j+\floor{(p^2-j)/p}\ge p+1.
\]
For $j=p^2$ this is $p^2\ge p+1$. Otherwise $v_p\binom{p^2}{j}=2-v_p(j)$. If $j=ap+b$ with $1\le b<p$, the expression is $p+1+a(p-1)+b$; if $j=ap$, $1\le a<p$, it is $p+1+a(p-1)$. Applying the second identity to the leading term proves
\[
 D_X^q=p^p u_XD_X+p^{p+1}U_X
\]
with $u_X$ a unit and $U_X$ integral. Raise the second identity of \eqref{eq:dual-low-ops} to $p$. Its leading term uses the first identity; every mixed term $1\le j<p$ has valuation at least $1+p(p-j)+(p+1)j=p^2+j+1$, and its last term has valuation $p(p+1)\ge p^2+1$. Hence
\[
 D_Y^q=p^{p^2}u_YD_Y+p^{p^2+1}U_Y.
\]

For either direction put $(r,B,c)=(p^2,p,c_X)$ or $(p,p^2,c_Y)$. Then $c=B/(q-1)<1$ and $cr=q/(q-1)$. For $0\le n<q$, the maximum of $cn-\floor{n/r}$ is one, attained at $n=q-1$: within a block it increases, and successive block maxima increase by $cr-1>0$. This proves the affine inequality at the initial orders. For $n\ge q$, apply the common recurrence $D^q=p^BuD+p^{B+1}U$ to $D^{n-q}f$. Induction bounds the first term by $cn-1$ and the second by $cn-1+(1-c)$. Integrality and integer valuations give the ceiling. Since the bounds are on the whole formal function ring, composing the two operator containments proves the mixed assertion.

The affine constants are attained by the same monomial $f=X^{p-1}Y^{p^2-1}$ at $n=q-1$. Along the $X$ logarithmic axis the leading terms of the exponential are $X=t+\cdots$, $Y=-t^p/p+\cdots$, so $f$ has first degree $q-1$ with coefficient $(-1/p)^{p^2-1}$. Since $v_p((q-1)!)=p^2+p-2$, the derivative has valuation $p-1=c_X(q-1)-1$. Along the $Y$ axis the leading terms are $X=-t^{p^2}/p+\cdots$, $Y=t+\cdots$; the coefficient is $(-1/p)^{p-1}$ and the derivative valuation is $p^2-1=c_Y(q-1)-1$.

For the additional exact sequences, DP32's scaled inverse here reduces to
\[
 \overline{\mathcal E'}_1(u,v)=\sum_{r\ge0}u^{q^r}-\sum_{r\ge0}v^{p^2q^r},\qquad
 \overline{\mathcal E'}_2(u,v)=\sum_{r\ge0}v^{q^r}-\sum_{r\ge0}u^{pq^r}.
\]
The nonzero off-diagonal coefficients, the respective input/output weights $r_X=(p^2+1)/(q-1)$ and $r_Y=(p+1)/(q-1)$, and Legendre's formula give exactly the two stated valuations. The diagonal sequences from DP32 prove slope optimality. In particular $D_X^pY(0)=-(p-1)!$ is a unit, while $c_X>0$.
\end{proof}
